%=======================================================================
%  Paper V — Section 8b / 9 (Practical guidelines)   FINAL
%
%  Requires the macros listed at the head of Section 2.
%
%  PRINCIPAL ADDITIONS IN THIS REVISION ------------------------------
%  (a) S is shown to equal Delta M / delta(Delta M) identically, so it
%      is the signal-to-noise ratio of the excess mass and is invariant
%      under exchanging the labels "region" and "complement". A and the
%      raw curvature interval are not. This resolves the concern raised
%      in Section 7.5 that S might be a coordinate artefact, and
%      explains why A fails there: A is label-dependent.
%  (b) The reporting rule is restated in terms of the LARGER piece.
%      From the mass constraint, |d rho_rest| / |d rho_R| = f/(1-f), so
%      the better-determined density is always that of the larger
%      piece, with the threshold exactly at f = 1/2. On Arrokoth the
%      region holds 66.3 %, so "report rho_rest" as previously worded
%      points at the worse-determined of the two by a factor 1.83.
%  (c) The three-part detection criterion of Section 4.4 and the
%      diagnostic-versus-failure table are brought in, since no single
%      number separates a result from an artefact.
%
%  Superseded statements:
%      "independent parameters of the same fit"
%              -> affinely related; one free parameter (Sections 2.2, 5.1)
%      "a fractional range in Delta M is SMALLER than in rho_rest by
%       rho_rest V_tot / Delta M"
%              -> LARGER, by that factor; Eq. (excessamplify)
%      "7.5 kg/m3 per metre in the linear regime"
%              -> 10.0 at the reference plane, 55.0 at the outer planes;
%                 and the quantity to quote is rho_rest, at 0.28
%      "place the smaller of the two pieces in the region" (Section 7.3)
%              -> relabelling buys no precision; it determines which
%                 density the reporting rule names
%      "the case for a real anomaly rests on Lowry et al."
%              -> that determination is 1.75 sigma from zero and its
%                 interval contains every plane of the sweep
%      "0.94 % of the surface area"  -> 0.943 %
%
%  All values confirmed. Resolved in this pass: the 1.35 span threshold is
%  replaced by reporting span and standard error, consistently with Section
%  5.2; the envelope row now carries the Eros-specific value, 6.03 % with
%  quadric decimation and 5.40 % with clustering, and the note records that
%  the entry turns into a cross under the second.
%=======================================================================

\section{Practical guidelines}
\label{sec:guidelines}

The results above are mostly negative or qualified, which makes them
awkward to act on. This section states what they imply for anyone applying
the method to a new body. Nothing here is new evidence with one exception,
noted where it arises: the relation between the displacement statistic and
the excess mass in Section~\ref{sec:guidelines:report} is an identity we
had not previously drawn out, and it changes which quantity we recommend
reporting. The rest is the preceding sections read as instructions.

\subsection{When the method is worth applying}
\label{sec:guidelines:when}

Four conditions should hold before the result will mean much.

\emph{The mass must be measured.} On a body without one, the assumed bulk
density can carry the recovered contrast to the threshold of a sign
reversal, as it does on Arrokoth within the plausible range
(Section~\ref{sec:nulls:arrokoth}), where the recovered contrast runs from
$0.9999$ at $\SI{250}{\kilo\gram\per\cubic\metre}$ to $0.9159$ at $500$ and
the plane sweep at fixed assumed mass carries $\Dmass$ from $+2.5\%$ to
$-11.6\%$ of the total. A result that depends on an assumption in that way
is not a measurement. Where no mass exists the method may still be run, but
the honest report is a range that includes both senses.

\emph{The body should have a dividing surface that something other than the
objective function can locate.} The dispersion does not constrain the plane
(Section~\ref{sec:plane}), so a saddle in the cross-sectional profile, or
some comparable geomorphological feature, has to supply it. Bennu, which has
no saddle, is the case where this fails: a plane can be imposed, but nothing
justifies its position.

\emph{The relaxation premise should be plausible for the particular body.}
Section~\ref{sec:discussion:premise} assesses it for the five studied here,
and finds it strong on Eros, weak on 67P and Arrokoth, uncertain on Bennu,
whose surface may still be in motion, and failing its first stated condition
on Itokawa. One qualification belongs with that last item, because it is
easily over-read. Section~\ref{sec:discussion:conditioning} shows that the
first condition---that the rotational term be a significant fraction of the
gravitational one---is the condition under which the \emph{bulk-density}
inversion of \citet{richardson2014} is not identically flat, and that it has
no counterpart in a fixed-mass two-region inversion, where the mass
constraint fixes the scale instead. A slow rotator is therefore not
ill-conditioned for the problem solved here. What remains true, and is the
reason the condition still appears in this list, is that the premise itself
is a statement about whether the surface has relaxed, and the empirical
evidence that it has is weakest where the rotation is slowest. A body
resurfaced by activity, or one whose surface is still in motion, is not a
good candidate whatever its spin.

\emph{The anticipated structure must be one a plane can represent.} A radial
or axisymmetric arrangement, of the kind \citet{scheeres2020} report for
Bennu, is invisible to this model however well the other conditions hold.

\subsection{What to run}
\label{sec:guidelines:what}

A single inversion at a single plane on a single mesh is not enough to know
whether its output means anything. The following are cheap, and each answers
a question the others cannot.

\emph{Vary the mesh resolution, and report both the density and the plane.}
Two or three meshes establish whether the recovered density is stable, which
the stability of the dispersion does not by itself guarantee
(Section~\ref{sec:itokawa:mesh}). Run this twice. Holding the plane fixed
isolates the sensitivity of the objective function; letting the detector
choose it on each mesh measures the sensitivity of the plane, and on Itokawa
the second is much the larger of the two---the detected plane moves by
$\SI{9.4}{\metre}$ between the finest and coarsest meshes, displacing
$\denshead$ by about $3.4\%$ against the $0.45\%$ to which the fixed-plane
comparison bounds it. Report the detected plane on every mesh, not the
density alone; a study that fixed the plane externally would not see this
error source at all.

\emph{Sweep the dividing plane, selecting by volume fraction.} Select planes
by volume fraction rather than by position, so that the sweep spans a
comparable range of region volume on bodies of different shape, and report
the scaling exponent of Eq.~\eqref{eq:slope} together with the fractional
ranges of $\densR$, $\densrest$ and $\Dmass$ and the subset of planes over
which each was computed. Three conditions govern whether the exponent means
anything. The density difference must keep its sign across the sweep, since
a quantity passing through zero has no fractional range and no logarithmic
slope; the sweep must span an adequate range in $\VR$, which we judge by
reporting the span and the standard error of the fit with every exponent
rather than by imposing a threshold---the spans used here are $1.907$ on
Itokawa, $1.504$ on 67P and $1.90$ on every synthetic body, all far above
the factor of about $1.35$ below which the logarithmic fit ceases to be
constrained; and the exponent should not be presented as
independent evidence about $\densrest$, since Eq.~\eqref{eq:steepidentity}
makes $s+1$ identically the logarithmic drift of $\densbulk-\densrest$. What
the exponent adds is the distinction between a systematic drift and scatter
about a fixed value, which a coefficient of variation cannot make, and a
well-defined quantity where fractional ranges diverge. Report the
centre-of-mass offset across the same sweep as well: comparing its scaling
against the purely geometric slope of Eq.~\eqref{eq:geomslope} is what
excludes conservation of the mass dipole
(Section~\ref{sec:constrains:dipole}), and it costs nothing, the offset
being a by-product of the same solutions.
% Resolved consistently with Section 5.2: the threshold does no work, since
% no sweep reported here comes near it, and the standard error is now quoted
% with every exponent (Itokawa -1.141 +/- 0.018, 67P -2.80 +/- 0.27, and the
% four synthetic exponents in Table 5.2).

\emph{Vary the assumed mass.} Three or four runs over a factor of a few show
whether the answer is driven by the data or by the assumption
(Section~\ref{sec:constrains:mass}). Read the outcome as the position of the
resulting contrast interval relative to unity and not as its width: on the
raw change in contrast Itokawa is the \emph{more} sensitive of the two
bodies tested, and it is only per e-fold of assumed mass, and in whether the
interval reaches unity, that the two separate. We put this forward as a
candidate diagnostic rather than an established one, since it rests on two
bodies here, and note the awkwardness that it is of use precisely where the
mass is unknown and there is nothing to calibrate it against.

\emph{Check that no minimum sits on a scan boundary, and widen the interval
until none does.} A minimisation reports the smallest value it finds, and a
minimum at the edge of the interval searched is the edge, not a result. This
guard fired twice in the present study: on the Arrokoth neck model, where a
slab of $1.0\%$ of the volume drove the density to the scan floor and
returned an improvement of $13.6\%$ that had to be discarded, and at one
Itokawa plane, where widening the interval returned an interior minimum
(Section~\ref{sec:failures:edge}). Record the interval used at every plane,
since it may have to change along a sweep, and check separately that the
upper limit is not implicitly restricting the region size: holding $\Dmass$
fixed, a ceiling $\densR^{\max}$ admits no region smaller than
$\Dmass/[(\densR^{\max}-\densrest)\Vtot]$, which on Itokawa is $2.56\%$ and
bounds the range of planes the configuration can resolve at all.

\emph{Check the geometry every time, in five layers.} The region must agree
with a closed-form solid where one exists, satisfy its defining inequality
at every vertex, close against its complement, form a closed surface with
the Euler characteristic its component count implies, and satisfy the mass
identity $\Dmass=(\densbulk-\densrest)\Vtot=\Ddens\VR$ on every solution.
The last of these is the one that operates on the output rather than on the
geometry, and it is the layer that would have caught the construction error
of Section~\ref{sec:methods:cone}, which passed a volume-closure test while
enclosing a solid wrong by a factor of $1.6$. Where the shape model has more
than one component, test facet centroids for containment in the others as
well: none of the other checks detects a facet that lies inside the body
(Sections~\ref{sec:methods:verify} and~\ref{sec:failures:buried}).

\emph{Where the signal is weak, test the field-point offset.} The inward
displacement of field points from the facet centroids shifts the surface
potential by of order $0.04\%$ of $\lvert\Ubar\rvert$ on the bodies here
(Section~\ref{sec:methods:model}), which is comparable to the dispersion
improvements returned by all four weak-signal bodies. Repeating the
inversion with the offset scaled by ten and by a tenth establishes whether a
sub-per-cent improvement is a property of the body or of the discretisation.

\subsection{What to report}
\label{sec:guidelines:report}

\emph{Report the density of the larger of the two pieces, and the excess
mass, not the region density as such.} This is a correction to the rule
stated in earlier drafts, which was to report the background density
$\densrest$, and the correction matters on one of the five bodies here.
Differentiating the mass constraint~\eqref{eq:massbalance} at fixed total
mass gives
%
\begin{equation}
\frac{\lvert\delta\densrest\rvert}{\lvert\delta\densR\rvert}
 = \frac{\VR}{\Vtot-\VR} = \frac{f}{1-f} ,
\label{eq:whichpiece}
\end{equation}
%
so the density of the larger piece is always the better determined, in
absolute terms, and the threshold is exactly $f = \tfrac{1}{2}$. On Itokawa,
Eros and 67P the region is the smaller piece and $\densrest$ is accordingly
the constrained quantity, which is the case the earlier wording had in mind.
On Arrokoth the detector places $66.3\%$ of the volume in the region, so
$\densrest$ is the density of the \emph{smaller} piece and is the worse
determined of the two by a factor of $1.83$ in fractional terms. Reporting
it there, as Table~\ref{tab:nulls} does, points at the wrong number. Which
piece is called the region is a labelling choice and carries no physics: it
does not change the uncertainty on either density, only which of them the
reporting rule names.

Report $\Dmass$ alongside, since Eq.~\eqref{eq:excessidentity} makes it the
same statement and it is the quantity that survives the relabelling. The
region density is a derived quantity depending on a dividing surface the
investigator chose (Section~\ref{sec:constrains:claim}), and a region
density quoted without its plane is not reproducible.

\emph{Quote the amplification factor with any stability claim, and the plane
at which it was evaluated.} A fractional range in $\densrest$ corresponds to
a fractional range in $\Dmass$ that is \emph{larger} by
$\ampfac = \densrest\Vtot/\Dmass$, not smaller---on Itokawa,
$\CV(\Dmass)=2.95\%$ against $\CV(\densrest)=0.251\%$, a ratio of $11.76$
which reproduces the mean $\ampfac$ of $11.77$ over the same planes. Any
comparison of how well two quantities are determined must divide by it
first. When comparing $\densR$ against $\densrest$, note that the two are
comparable because they are densities of the same model in the same units,
and \emph{not} because they are independent: the mass constraint leaves
exactly one free parameter and relates them affinely as
$\densrest = a - b\,\densR$ (Section~\ref{sec:methods:identity}). Quote
$\ampfac$ with its plane, since Eq.~\eqref{eq:ampsensitivity} makes it move
by $\pm5.6\%$ across a sweep over which $\densrest$ moves by $\pm0.434\%$,
and note that a large $\ampfac$ is a warning rather than a precision: a body
returning $\ampfac\simeq1200$, as Bennu does, is one on which the model has
found essentially no excess mass.

\emph{State the component count of the shape model, and whether any facets
were excluded.} A model delivered as several unjoined shells may have part
of one inside another; on Arrokoth this affected $0.943\%$ of the surface
area and changed a qualitative conclusion
(Section~\ref{sec:nulls:arrokoth}).

\emph{State the plane, the criterion that placed it, and the sensitivity of
the reported quantity to moving it.} The sensitivity must be quoted for the
quantity being reported, and the three differ by two orders of magnitude. At
the Itokawa reference plane they are
$\SI{10.0}{\kilo\gram\per\cubic\metre\per\metre}$ for $\denshead$---rising
to $\SI{55}{\kilo\gram\per\cubic\metre\per\metre}$ at the outer planes and
falling to $7.5$ over the first four, so no single figure describes the
sweep---$\SI{0.28}{\kilo\gram\per\cubic\metre\per\metre}$ for $\densrest$,
and $0.18\%$ per metre for $\Dmass$. Earlier drafts quoted the first alone,
which states the sensitivity of the quantity
Section~\ref{sec:constrains} concludes the method does not determine.

\emph{State the saddle prominence, and do not read it as a figure of merit.}
It lets a reader judge the provenance of the plane, and the values here
separate the accepted from the refused by more than two orders of magnitude,
from $0.091$ on 67P down to $0.0004$ on Bennu. Among the accepted bodies it
carries no information about the strength of the result: Arrokoth has the
highest prominence in the study, $0.490$, and returns an improvement of
$1.272\%$, while Itokawa has the third, $0.170$, and returns $16.9\%$.

\emph{Give the facet count used, not the one requested, and the decimation
method.} A slope or dispersion statistic depends on both
(Section~\ref{sec:discussion:ellipsoid}).

\emph{Bound the answer by the material, and state that the bound is
one-sided.} A recovered density above the grain density of the meteorite
analogue is not a result but a signal that the plane has advanced too far:
on Itokawa the outermost plane of the sweep exceeds the LL-chondrite grain
density by $11.2$ standard deviations
(Sections~\ref{sec:itokawa:ceiling} and~\ref{sec:plane:impossible}). Convert
the bound into a plane position through the measured gradient and carry its
uncertainty: on Itokawa the block-density bound of
$3220\pm\SI{220}{\kilo\gram\per\cubic\metre}$ becomes
$\xsplit \le \SI{194+-11}{\metre}$, and quoting it to the tenth of a metre
would overstate it. The bound excludes planes that are too far out and says
nothing about planes that are too far in, so it truncates a sweep and does
not locate one.

\subsection{How to read the outcome}
\label{sec:guidelines:read}

\emph{Require three things before calling a result a detection.} No single
number separates a result from an artefact, which is the lesson of
Section~\ref{sec:failures} and the reason for stating a criterion rather
than a threshold. The improvement must exceed the value the mesh-resolution
envelope takes on that body, not the upper edge of a range attained on some
other body; the displacement statistic $\Sstat$ must be large; and the
dividing plane must carry geometric warrant, meaning a saddle returned by a
detector at a prominence above threshold rather than a plane chosen by the
investigator. The misplaced plane of Section~\ref{sec:failures:plane}
satisfies the second and fails the first and third. Itokawa satisfies all
three.

\emph{Use $\Sstat$ rather than a curvature interval, because it is the only
one of the available diagnostics that does not depend on which piece was
called the region.} Its definition in Eq.~\eqref{eq:sstat} is a displacement
measured in units of an uncertainty, but multiplying numerator and
denominator by $\Vtot$ turns it into something more directly meaningful:
%
\begin{equation}
\Sstat = \frac{\lvert\densbulk-\densrest\rvert}{\sigma(\densrest)}
       = \frac{\lvert\densbulk-\densrest\rvert\,\Vtot}
              {\sigma(\densrest)\,\Vtot}
       = \frac{\lvert\Dmass\rvert}{\delta(\Dmass)} ,
\label{eq:ssnr}
\end{equation}
%
which is exact. $\Sstat$ is the signal-to-noise ratio of the excess mass,
and since $\Dmass$ merely changes sign when the labels ``region'' and
``complement'' are exchanged, $\Sstat$ is invariant under that exchange. The
raw curvature interval is not, and neither is $\ampfac$: at the misplaced
Eros plane the interval is $\pm15$ under one labelling and $\pm115$ under
the other, and $\ampfac$ is $-9.12$ under one and $+61.0$ under the other.
That is why $\ampfac$ ranks the misplaced plane as the strongest signal in
the study, and it is a limitation of the quantity rather than a statistical
accident: $\ampfac$ is a cheap first check on bodies whose planes are
properly placed (Section~\ref{sec:nulls:amplification}) and says nothing
about whether a plane is properly placed. Read the other way,
Eq.~\eqref{eq:ssnr} says that $\Sstat = 5.8$ on Itokawa is a statement that
the excess mass is determined to $1/5.8 = 17\%$, which is the figure the
error budget of Eq.~\eqref{eq:budget} arrives at independently.

\emph{A strong response is not by itself evidence of a physical anomaly.}
The synthetic bodies of Section~\ref{sec:constrains:synthetic} are uniform
by construction and still produce strong responses and a pinned background
density when their shape is sufficiently asymmetric. What that pinning
identifies is a regime of the objective function. On Itokawa the case for a
real anomaly rests on the independent YORP determination of
\citet{lowry2014}, not on the inversion alone---but that support should not
be asked to carry more than it can. The Lowry offset of
$\SI{21+-12}{\metre}$ is itself only $1.75$ standard deviations from zero,
so it does not exclude a homogeneous Itokawa at two sigma, and its interval
contains the centre-of-mass offset returned at \emph{every} plane of
Table~\ref{tab:itokawasweep}. It corroborates the sign and the order of
magnitude of the anomaly and discriminates neither the plane nor the value.

\emph{A weak response is not evidence of homogeneity.} It is consistent with
a uniform interior, with a structure the planar model cannot see, and with
the relaxation premise failing on that body. These are not distinguishable
from within the method (Section~\ref{sec:nulls}).

\emph{A sharp minimum is not a precise answer; propagate the interval before
reading it.} The tightest curvature bound in this study,
$\pm\SI{15}{\kilo\gram\per\cubic\metre}$, belongs to the one run known to be
wrong: a plane placed on a flank of Eros, where the complement holds only
$11.56\%$ of the volume and the mass constraint amplifies the objective on
both sides of the minimum. Carried through Eq.~\eqref{eq:whichpiece} that
interval becomes $\SI{115}{\kilo\gram\per\cubic\metre}$ on the complement
density, or $3.83\%$---the worst determination anywhere in this study, and
correspondingly $\Sstat = 2.86$ against $5.8$ on Itokawa. The apparent
precision is a property of the coordinate in which the scan was performed.

\emph{Watch for an improvement bought with almost no mass.} The ratio of the
dispersion improvement to the fractional excess mass is $14.4$ for the
Arrokoth neck model, which the boundary guard discarded, against $2.20$ for
the Itokawa head and below unity for every other case in the study
(Section~\ref{sec:failures:edge}). It separates an edge minimum from a
genuine result by a factor of $6.5$ and costs nothing to compute. Its
limitation should be stated with it: it does not flag a misplaced plane,
whose ratio is unremarkable.

\emph{Corroboration should be sought outside the method.} Itokawa is
interpretable here because a YORP measurement, resting on different
premises, gives a compatible centre-of-mass offset, and because
Section~\ref{sec:constrains:dipole} establishes that the objective function
demonstrably does not hold the mass dipole fixed, so the two lines of
evidence are independent rather than two readings of one. Where no such
check exists, the result should be reported as what the objective function
prefers, not as what the interior is.

\begin{table}[htbp]
\centering
\scriptsize
\setlength{\tabcolsep}{3pt}
\caption{Which diagnostic responds to which failure, from the cases in
Section~\ref{sec:failures}. A tick marks a diagnostic that identifies the
failure, a cross one that is silent or actively misleading, and a dash one
that does not bear on it. The table is the argument for a three-part
criterion rather than a threshold: no row catches everything, and two
diagnostics rank the misplaced plane \emph{above} the genuine result. The
last two columns are the two silent failures that no statistical diagnostic
reaches, and that only geometric checks or the provenance of the plane will
find.}
\label{tab:diagnostics}
\setlength{\tabcolsep}{0.2pt}
\begin{tabular}{lccccc}
\toprule
Diagnostic & Region too small & Complement too small & Edge minimum
 & Misplaced plane & Wrong solid \\
\midrule
Curvature interval in $\densR$
 & \checkmark\ $+22.3\%$ & $\times$\ $0.57\%$ & --- & $\times$\ $0.57\%$ & --- \\
Interval propagated to the larger piece
 & $\times$\ $1.01\%$ & \checkmark\ $3.83\%$ & --- & \checkmark\ $3.83\%$ & --- \\
Boundary-minimum guard
 & --- & --- & \checkmark & --- & --- \\
$\improve$ per unit $\lvert\Dmass\rvert/\Mtot$
 & --- & --- & \checkmark\ $14.4$ & $\times$\ $0.45$ & --- \\
$\Sstat$
 & --- & \checkmark & --- & partly, $2.86$ & --- \\
$\ampfac$
 & --- & --- & --- & $\times$\ $9.12$ & --- \\
Envelope comparison$^{a}$
 & --- & --- & --- & narrowly, $5.53\%$ & --- \\
Plane provenance
 & --- & --- & --- & \checkmark & --- \\
Vertex predicate; mass-balance audit
 & --- & --- & --- & --- & \checkmark \\
\bottomrule
\end{tabular}
\end{table}
\vspace{2pt}
\begin{minipage}{\textwidth}
\footnotesize
\textsuperscript{a}\,The Eros-specific envelope is $6.03\%$ with quadric
decimation, so the $5.53\%$ returned at the misplaced plane is $0.92$ times
it and the comparison rejects the case. With vertex clustering the envelope
is $5.40\%$, the same result is $1.02$ times it, and the entry becomes a
cross. The margin is therefore smaller than the difference between two
defensible ways of computing the envelope, which is why the row reads
``narrowly'' and why plane provenance, the row below it, is the check we
recommend relying on.
\end{minipage}
